\section{模型评价}

\subsection{模型优点}

\begin{enumerate}[label=（\arabic*）]
\item GAMM处理了纵向数据的面板结构（ICC=0.60），通过随机截距和随机斜率建模个体异质性，以B样条捕捉孕周对游离DNA浓度的非线性效应。似然比检验（$\chi^2=76.76$，$p<1\times10^{-15}$）证实非线性模型优于线性混合模型，AIC改善66.8个单位。

\item 双风险优化框架将NIPT检测时点选择转化为决策理论问题。风险函数同时纳入延迟风险（sigmoid衰减）和浓度不足概率（基于GAMM预测分布），为不同BMI组给出可量化的最优时点，避免了单纯依赖经验阈值的局限。

\item 残差分解策略解决了BMI与体重之间的高度共线性。提取体重对BMI和身高回归后的残差作为独立变量，VIF从351降至1.01，保留了全部体成分信息的同时消除了参数估计的多重共线性膨胀。

\item 蒙特卡洛误差传播（500次模拟）为每个分组的最优时点提供了标准差和置信区间，诚实地暴露了高BMI组（G3标准差3.04周、G4标准差4.75周）的统计不确定性，避免给出表面精确但实际不可靠的推荐。

\item 马氏距离异常检测利用6个特征之间的协方差结构，在多维空间中度量样本偏离程度。相比逐变量Z值方法（AUC=0.492），马氏距离方法（AUC=0.574）能识别单一变量无法捕捉的联合异常模式。
\end{enumerate}

\subsection{模型不足}

\begin{enumerate}[label=（\arabic*）]
\item 样本规模有限，男胎数据仅267名孕妇，高BMI组（BMI$>$36）仅20人。蒙特卡洛分析显示该组最优时点的95\%置信区间跨度超过13周，时点推荐对这一人群缺乏实际指导价值。BMI$>$40的极端组仅4例，无法支撑统计推断。

\item GAMM残差峰度为4.15，高于正态分布的3.0，存在未被模型捕捉的厚尾变异。Beta-GLMM理论上更匹配Y浓度的$(0,1)$有界分布特征，但因收敛困难未能实现。这一分布假设偏差可能影响尾部预测的准确性。

\item 问题4的马氏距离模型AUC仅为0.574，略高于随机水平。最佳F1=0.217，召回率26.9\%，精确率18.2\%。异常与正常样本的特征分布高度重叠，6维特征空间无法提供足够的分离度。根本原因在于游离DNA浓度与染色体异常之间的生物学关联较弱，并非方法选择的问题。

\item 风险函数中$w_d=0.4$和$w_a=0.6$基于合理假设设定，缺乏临床结局数据的实证校准。灵敏度分析表明高BMI组的最优时点对权重选择高度敏感（$w_d$从0.2变至0.6时G3组偏移4.1周），权重的不确定性直接传递至推荐结果。

\item 受限于面板数据的小样本结构，未执行交叉验证。留一患者交叉验证（LOPOCV）是混合模型框架下评估泛化能力的理想方案，但每折需重新拟合GAMM，计算开销较大，在当前数据规模和模型复杂度下未能实施。
\end{enumerate}
